\section{Integral in Coordinates}

Firstly, let's represent the integral kernel $\mathcal{K}$ in vector form:
\begin{equation*}
\mathcal{K}= \mathcal{R}\left[ e^{*}_{\alpha \beta}(\hat{\Omega}) \frac{1}{2} \frac{\hat{p}^{\alpha} \hat{p}^{\beta}}{1 + \hat{\Omega} \cdot \hat{p}} \cdot e_{\mu \nu}(\hat{\Omega}) \frac{1}{2} \frac{\hat{q}^{\mu} \hat{q}^{\nu}}{1 + \hat{\Omega} \cdot \hat{q}}
\right],
\end{equation*}
\begin{equation*}
\mathcal{K}= \frac{1}{4} \frac{\hat{p}^{\alpha} \hat{p}^{\beta}\hat{q}^{\mu} \hat{q}^{\nu}}{(1 + \hat{\Omega} \cdot \hat{p})(1 + \hat{\Omega} \cdot \hat{q})} \mathcal{R}\left[e^{*}_{\alpha \beta}(\hat{\Omega}) \cdot  e_{\mu \nu}(\hat{\Omega})\right],
\end{equation*}
where:
\begin{equation*}
\mathcal{R}\left[e^{*}_{\alpha \beta}(\hat{\Omega}) \cdot  e_{\mu \nu}(\hat{\Omega})\right] = e^{+}_{\alpha \beta}(\hat{\Omega})e^{+}_{\mu \nu}(\hat{\Omega}) + e^{\times}_{\alpha \beta}(\hat{\Omega})e^{\times}_{\mu \nu}(\hat{\Omega}),
\end{equation*}
then:
\begin{equation}
\mathcal{K} = \frac{1}{4} \frac{\hat{p}^{\alpha} \hat{p}^{\beta}\hat{q}^{\mu} \hat{q}^{\nu} \left[ e^{+}_{\alpha \beta}(\hat{\Omega})e^{+}_{\mu \nu}(\hat{\Omega}) + e^{\times}_{\alpha \beta}(\hat{\Omega})e^{\times}_{\mu \nu}(\hat{\Omega}) \right]}{(1 + \hat{\Omega} \cdot \hat{p})(1 + \hat{\Omega} \cdot \hat{q})}.
\end{equation}

Consider individual tensor convolutions:
\begin{subequations}
\begin{align}
\hat{v}^{\alpha} \hat{v}^{\beta} e^{+}_{\alpha \beta}(\hat{\Omega}) &= (\hat{v} \cdot \hat{m})^2 - (\hat{v} \cdot \hat{n})^2, \\
\hat{v}^{\alpha} \hat{v}^{\beta} e^{\times}_{\alpha \beta}(\hat{\Omega}) &= 2 (\hat{v} \cdot \hat{m})(\hat{v} \cdot \hat{n}).
\end{align}
\end{subequations}
Let's denote $P_{a} = (\hat{p} \cdot \hat{a})$ and $Q_{a} = (\hat{q} \cdot \hat{a})$, then:
\begin{equation}
\mathcal{K} = \frac{1}{4} \frac{(P_m^2 - P_n^2)\cdot(Q_m^2 - Q_n^2) + 4 P_m P_n Q_m Q_n}{(1 + P_{\Omega})(1 + Q_{\Omega})}.
\end{equation}

Note that since $(\hat{m}, \hat{n}, \hat{\Omega})$ is an orthonormal basis, and $|\hat{p}| = 1$, $|\hat{q}| = 1$, it means that:
\begin{align*}
    P_n^2 &= 1 - P_m^2 - P_{\Omega}^2, \\
    Q_n^2 &= 1 - Q_m^2 - Q_{\Omega}^2.
\end{align*}

Then:
\begin{align*}
\frac{(P_n^2 - P_m^2)}{1 + P_{\Omega}} &= \frac{1 - 2 P_m^2 - P_{\Omega}^2}{1 + P_{\Omega}} = (1 - P_{\Omega}) - \frac{2 P_m^2}{1 + P_{\Omega}}, \\
\frac{(Q_n^2 - Q_m^2)}{1 + Q_{\Omega}} &= \frac{1 - 2 Q_m^2 - Q_{\Omega}^2}{1 + Q_{\Omega}} = (1 - Q_{\Omega}) - \frac{2Q_m^2}{1 + Q_{\Omega}}.
\end{align*}

As a result, we can divide the integral kernel into four parts:
\begin{equation}
\mathcal{K}(\hat{\Omega},\hat{p}, \hat{q}) = \frac{1}{4}(\mathcal{K}_{I} + \mathcal{K}_{J} + \mathcal{K}_{K} + \mathcal{K}_{H}),
\end{equation}

where:
\begin{subequations}
\begin{align}
\mathcal{K}_{I} &= (1 - P_{\Omega})(1 - Q_{\Omega}), \\
\mathcal{K}_{J} &=  - 2 \cdot \frac{P_m^2 (1 - Q_{\Omega})}{(1+P_{\Omega})}, \\
\mathcal{K}_{K} &=  - 2 \cdot \frac{Q_m^2 (1 - P_{\Omega})}{(1+Q_{\Omega})}, \\
\mathcal{K}_{H} &= 4 \cdot \frac{P_m Q_m (P_m Q_m + P_n Q_n)}{(1+P_{\Omega})(1+Q_{\Omega})},
\end{align}
\end{subequations}
and we can divide the integral into the same four parts:
\begin{equation}
\Gamma^{(g)}(\hat{p}, \hat{q}) = \frac{3 \cdot G (\kappa)}{4\pi} \cdot \frac{1}{4} \left( I + J + K + H \right),
\end{equation}

where:
\begin{subequations}
\begin{align}
I &= \int_{S^2} d\hat{\Omega} \cdot \mathcal{K}_{I} \cdot e^{\kappa \hat{\Omega} \cdot \hat{\Omega}_0} \\
J &= \int_{S^2} d\hat{\Omega} \cdot \mathcal{K}_{J} \cdot e^{\kappa \hat{\Omega} \cdot \hat{\Omega}_0} \\
K &= \int_{S^2} d\hat{\Omega} \cdot \mathcal{K}_{K} \cdot e^{\kappa \hat{\Omega} \cdot \hat{\Omega}_0} \\
H &= \int_{S^2} d\hat{\Omega} \cdot \mathcal{K}_{H} \cdot e^{\kappa \hat{\Omega} \cdot \hat{\Omega}_0}
\end{align}
\end{subequations}

Let's define a convenient coordinate system. As we can see, all integrals include a term containing $(\hat{\Omega} \cdot \hat{\Omega}_0)$ in the exponent, and to simplify this term, we take $\hat{\Omega}_0 \equiv (0, 0, 1)$, because then $(\hat{\Omega} \cdot \hat{\Omega}_0) = \cos \theta$. For the pulsar vectors, we have:
\begin{subequations}
\begin{align}
    \hat{p} &= (\sin \alpha, 0, \cos \alpha) \\
    \hat{q} &= (\sin \beta \cos \gamma, \sin \beta \sin \gamma, \cos \beta)
\end{align}
\end{subequations}
where $\alpha$, $\beta$ are the angles between the source center and the pulsars, and $\gamma$ is the angle between the directions to the pulsars from the source center. Now we can calculate all the scalar products:
\begin{align*}
    P_{\Omega} &= \sin \alpha \sin \theta \cos \phi + \cos \alpha \cos \theta, \\
    Q_{\Omega} &= \sin \beta \sin \theta \cos (\phi - \gamma) + \cos \beta \cos \theta, \\
    P_{n} &= \sin \alpha \cos \theta \cos \phi - \cos \alpha \sin \theta, \\
    Q_{n} &= \sin \beta \cos \theta \cos (\phi - \gamma) - \cos \beta \sin \theta \\
    P_{m} &= \sin \alpha \sin \phi, \\
    Q_{m} &= \sin \beta \sin (\phi - \gamma);
\end{align*}
and define scalar product between $\hat{p}$ and $\hat{q}$ as:

\begin{equation}
    \cos \xi = \sin \alpha \sin \beta \cos \gamma + \cos \alpha \cos \beta,
\end{equation}
where $\xi$ is the angle between $\hat{p}$ and $\hat{q}$.

Also, let's highlight the dependence of the coefficients on $\phi$:
\begin{subequations}
\begin{align}
    P_{\Omega} &= a \cos \phi + b, Q_{\Omega} = c \cos (\phi - \gamma) + d, \\
    P_{n} &= h \cos \phi - s, Q_{n} = k \cos (\phi - \gamma) - t, \\
    P_{m} &= f \sin \phi, Q_{m} = g \sin (\phi - \gamma);
\end{align}
\end{subequations}
where:
\begin{subequations}
\begin{align}
    a &= \sin \alpha \sin \theta, b = \cos \alpha \cos \theta, \\
    c &= \sin \beta \sin \theta, d = \cos \beta \cos \theta, \\
    h &= \sin \alpha \cos \theta, s = \cos \alpha \sin \theta, \\
    k &= \sin \beta \cos \theta, t = \cos \beta \sin \theta, \\
    f &= \sin \alpha , g = \sin \beta,
\end{align}

\end{subequations}